\section{Conclusion}

Diagonal isoduality turns the transversal-group problem into a one-dimensional
linear test.  MDS distance forces every nonzero code-to-dual diagonal
multiplier to be invertible and forces all such multipliers onto one projective
line.  That line either constructs the elementary unipotents generating
\(\mathrm{SL}_2(q)\), or its absence leaves only the split torus.  Stabilizer
AME rigidity then proves that these Clifford constructions exhaust every
product-unitary logical implementation.

For six coordinates, the same dichotomy becomes the conic criterion in the
six-arc model.  The pencil quotient, Clebsch syndrome geometry, and computed
party extensions show how the all-length criterion organizes finer finite
geometry without depending on those calculations.  Fixed-copy scalar
contractions are generically constant, whereas the operator-valued reduced
states retain the local Weyl directions.  This explains why the rigidity
argument succeeds where bounded-degree scalar invariants do not.

Two exact questions remain.  Over extension fields, the full additive
Clifford group can contain nonsemilinear sectors not detected by the prime-field
statement.  For party-moving symmetries, the nonabelian factor set gives an
exact splitting criterion, but a conceptual classification beyond the twelve
computed rows is open.
